\documentclass[11pt]{amsart}
\usepackage{geometry}
\geometry{paperwidth=210mm,paperheight=297mm,top=35mm,bottom=35mm,left=25mm,right=25mm}
\usepackage[colorlinks,pagebackref,hypertexnames=false]{hyperref}
\usepackage{microtype}
\usepackage{xcolor}
\usepackage{amsmath}
\usepackage{amstext}
\usepackage{amsfonts}
\usepackage{amssymb}
\usepackage{amsbsy}
\usepackage{latexsym}
\usepackage[T1]{fontenc}
\usepackage[all]{xy}
\usepackage{tikz-cd}
\usepackage{hhline}
\usepackage{enumerate}
\usepackage{enumitem}
\usepackage[nameinlink,noabbrev]{cleveref}
\usepackage{autonum}
\usepackage{dsfont}
\usepackage{booktabs} % For professional tables
\usepackage{caption} % Required for \captionof
\usepackage{thmtools}
\allowdisplaybreaks%%%% maybe will help

% Hyperref setup (you can keep or adjust colors)
\hypersetup{
	colorlinks=true,
	linkcolor=red,
	citecolor=blue,
	filecolor=blue,
	urlcolor=blue
}

\newenvironment{blueenv}
{\color{blue}} % Start coloring text blue
{}              % The color automatically ends with the environment


\newcommand{\vtx}[1]{*+[o][F-]{\scriptscriptstyle #1}} %the command to draw a vertex
\vfuzz2pt % Don't report over-full v-boxes if over-edge is small
\hfuzz2pt % Don't report over-full h-boxes if over-edge is small
\mathsurround=0pt \textwidth=15true cm \tolerance=1500%

\newcounter{num}[section] %
\renewcommand{\thenum}{\thesection.\arabic{num}} %


\newtheorem{mainthm}{Theorem}
\renewcommand{\themainthm}{\Alph{mainthm}}
\newtheorem{maincor}[mainthm]{Corollary}
\newtheorem{theorem}{Theorem}[section]
\newtheorem{Mtheorem}{Theorem}
\renewcommand*{\theMtheorem}{\Alph{Mtheorem}}
\newtheorem{Ctheorem}[Mtheorem]{(Theorem)}
\renewcommand*{\theCtheorem}{\Alph{Ctheorem}}


% --- THEOREMS, PROPOSITIONS, COROLLARIES, LEMMAS ---
\newenvironment{theo}
{\refstepcounter{num}%
	\bigskip\noindent{\bf Theorem~\arabic{section}.\arabic{num}.}\quad\itshape}
{\par\bigskip}

\newenvironment{prop}
{\refstepcounter{num}%
	\bigskip\noindent{\bf Proposition~\arabic{section}.\arabic{num}.}\quad\itshape}
{\par\bigskip}

\newenvironment{cor}
{\refstepcounter{num}%
	\bigskip\noindent{\bf Corollary~\arabic{section}.\arabic{num}.}\quad\itshape}
{\par\bigskip}

\newenvironment{lemma}
{\refstepcounter{num}%
	\bigskip\noindent{\bf Lemma~\arabic{section}.\arabic{num}.}\quad\itshape}
{\par\bigskip}

% --- EXAMPLES, REMARKS, CONJECTURES ---
\newenvironment{example}
{\refstepcounter{num}%
	\bigskip\noindent{\bf Example~\arabic{section}.\arabic{num}.}\quad}
{\par\bigskip} 

\newenvironment{remark}
{\refstepcounter{num}%
	\bigskip\noindent{\bf Remark~\arabic{section}.\arabic{num}.}\quad}
{\par\bigskip}

\newenvironment{conj}
{\refstepcounter{num}%
	\bigskip\noindent{\bf Conjecture~\arabic{section}.\arabic{num}.}\quad\itshape}
{\par\bigskip}

% --- DEFINITIONS ---
\newenvironment{definition}[1]
{\refstepcounter{num}%
	\bigskip\noindent{\bf Definition~\arabic{section}.\arabic{num}}~(\emph{#1}).\quad}
{\par\bigskip}

\newenvironment{defin}
{\refstepcounter{num}%
	\bigskip\noindent{\bf Definition~\arabic{section}.\arabic{num}.}\quad}
{\par\bigskip}

% --- NAMED ENVIRONMENTS ---
\newenvironment{prop_with_name}[1]
{\refstepcounter{num}%
	\bigskip\noindent{\bf Proposition~\arabic{section}.\arabic{num}} (#1).\quad\itshape}
{\par\bigskip}

\newenvironment{theo_with_name}[1]
{\refstepcounter{num}%
	\bigskip\noindent{\bf Theorem~\arabic{section}.\arabic{num}} (#1).\quad\itshape}
{\par\bigskip}

\newenvironment{remark_with_name}[1]
{\refstepcounter{num}%
	\bigskip\noindent{\bf Remark~\arabic{section}.\arabic{num}} (#1).\quad\itshape}
{\par\bigskip}

\newenvironment{lemma_with_name}[1]
{\refstepcounter{num}%
	\bigskip\noindent{\bf Lemma~\arabic{section}.\arabic{num}} (#1).\quad\itshape}
{\par\bigskip}

%\newenvironment{proof}{\noindent\emph{Proof. }}
%{$\Box$ \bigskip}
\newenvironment{proof_of}[1]{\noindent\emph{Proof #1}}
{$\Box$ \bigskip}
\newenvironment{eq}{\begin{equation}}{\end{equation}}

\renewcommand{\Ref}[1]{(\ref{#1})}

\newcommand{\si}{\sigma}
\newcommand{\al}{\alpha}
\newcommand{\be}{\beta}
\newcommand{\ga}{\gamma}
\newcommand{\la}{\lambda}
%\newcommand{\La}{\wedge}
\newcommand{\de}{\delta}
\newcommand{\De}{\Delta}
\newcommand{\La}{\Lambda}
\newcommand{\Ga}{\Gamma}

\newcommand{\LA}{\langle}
\newcommand{\RA}{\rangle}
\newcommand{\ov}[1]{\overline{#1}}
\newcommand{\un}[1]{{\underline{#1}} }
\newcommand{\id}[1]{{{\rm id}\{{#1}\}}}
\newcommand{\lin}{\mathop{\rm lin}}
\newcommand{\tr}{\mathop{\rm tr}}
\newcommand{\Ind}{\mathop{\rm Ind}}
%\newcommand{\sys}{{\mathcal{S}} }
%\newcommand{\projlim}{\mathop{\rm{projlim }}}
%\newcommand{\alg}{\mathop{\rm alg}}
\newcommand{\mdeg}{\mathop{\rm mdeg}}
\newcommand{\hterm}{\mathop{\rm ht}}   % the highest term 
\newcommand{\Skew}{\mathop{\rm skew}}   
\newcommand{\diag}{\mathop{\rm diag}}
%\newcommand{\mydiag}{\mathop{\rm d}}
\newcommand{\Char}{\mathop{\rm char}}
\newcommand{\image}{\mathop{\rm Im}}
\newcommand{\sign}{\mathop{\rm{sgn }}}
%\newcommand{\sgn}{\mathop{\rm{sgn }}}
\newcommand{\rank}{\mathop{\rm{rank }}}
%\newcommand{\Spec}{\mathop{\rm{Spec }}}
\newcommand{\trdeg}{\mathop{\rm{tr.deg }}}
%\newcommand{\St}{\mathop{\rm{St }}}
%\newcommand{\Orb}{\mathop{\rm{O }}}
\newcommand{\Hom}{{\mathop{\rm{Hom }}}}
\newcommand{\End}{{\mathop{\rm{End }}}}
\newcommand{\Ker}{{\mathop{\rm{Ker }}}}
\newcommand{\Balpha}{{\boldsymbol{\alpha}}}
\renewcommand{\d}{\mathrm{d}}
\newcommand{\dL}{\d^{\Lambda}}
\newcommand{\Om}{\Omega}
\newcommand{\C}{\mathbb{C}}
\newcommand{\X}{\breve{X}}
\DeclareMathOperator{\im}{Im}
\DeclareMathOperator{\gr}{gr}
% The two key operators
\newcommand{\dduL}{(\d + u\dL)}
\newcommand{\ddL}{(\d\dL)}




\newcommand{\Sym}{{\mathcal S}}
\newcommand{\wz}{\mathop{\rm wz}}
\newcommand{\upto}{,\ldots ,}
\newcommand{\st}{\,|\,}

%\DeclareMathOperator{\GL}{GL}
%\DeclareMathOperator{\SL}{SL}




\newcommand{\algA}{\mathcal{A}}    %algebra A
\newcommand{\algB}{\mathcal{B}}    %algebra B
\newcommand{\hsp}{Q}               % h.s.p.
%\newcommand{\X}{\LA X\RA}          % h.s.p.
\newcommand{\GL}{{\rm GL}}         % GL
\newcommand{\SL}{{\rm SL}}         % SL
\renewcommand{\O}{{\rm O}}         % GL
\newcommand{\SO}{{\rm SO}}         % SO
\newcommand{\Sp}{{\rm Sp}}         % Sp
\newcommand{\Op}{{\rm O}_2^{+}}    % O_2^+
%\newcommand{\Om}{{\rm O}_2^{-}}    % O_2^-

%\newcommand{\W}{\mathcal{W}} % the set of monomials in the generic matrices 
%\newcommand{\N}{\mathcal{N}} 
%\newcommand{\M}{\mathcal{M}}
%\newcommand{\T}{{\mathcal T}} % a tableau
%\renewcommand{\P}{{\mathcal P}} % 

\newcommand{\calV}{{\mathcal{V}}}
\newcommand{\calN}{{\mathcal{N}}}
\newcommand{\calU}{{\mathcal{U}}}
\newcommand{\calB}{{\mathcal{B}}}
\newcommand{\calD}{{\mathcal{D}}}
\newcommand{\calT}{{\mathcal{T}}}
\newcommand{\Hdr}{H_{dR}}
\newcommand{\Hdl}{H_{\dL}}
\newcommand{\Hharm}{\mathcal{H}} % Harmonic forms
\newcommand{\Lop}{L} % Lefschetz operator
\newcommand{\Lam}{\Lambda} % Adjoint






\newcommand{\FF}{{\mathbb{F}}}   % base field
\newcommand{\KK}{{\mathbb{K}}}   % some field
\newcommand{\CC}{{\mathbb{C}}}   % base field
\newcommand{\NN}{{\mathbb{N}}}
\newcommand{\ZZ}{{\mathbb{Z}}}   % integers
\newcommand{\QQ}{{\mathbb{Q}}}
\newcommand{\RR}{{\mathbb{R}}}

\renewcommand{\k}{\kappa}



\newcommand{\ii}{{\mathbb{I}}}   % (-1)^{1/2}

\newcommand{\Pf}{{\mathop{\rm{pf }}}}
\newcommand{\DP}{{\rm DP} }
%\newcommand{\F}{{\mathop{\rm{bpf }}}}
\newcommand{\bpf}{{\mathop{\rm{bpf }}}}

%\newcommand{\ovphi}[1]{\varphi(#1)}  

\newcommand{\QS}{\mathfrak{Q}} %Mixed quiver setting

\newcommand{\Q}{\mathcal{Q}}    %Quiver
\newcommand{\G}{\mathcal{G}}    %Quiver
\newcommand{\n}{\boldsymbol{n}} %dimension vector
\newcommand{\g}{\boldsymbol{g}} %for groups
\newcommand{\h}{\boldsymbol{h}} %for homomorphisms
\renewcommand{\i}{\boldsymbol{i}} %involution

\newcommand{\I}{\mathbb{I}} % i^2=-1

\newcommand{\matr}[4]{\left(
	\begin{array}{cc}
		#1 & #2 \\ 
		#3 & #4 \\ 
	\end{array}
	\right)}

\newcommand{\vect}[2]{
	\begin{pmatrix} #1\\ #2 \end{pmatrix}}

\newcommand{\besep}{\be_{\rm sep}}
\newcommand{\sisep}{\si_{\rm sep}}

\newcommand{\stabp}{\mathrm{Stab}^{+}\!}

%\newcommand{\e}{\mathfrak{e}}
\newcommand{\e}{\mathsf{e}}
%\newcommand{\e}{e}
\newcommand\restr[2]{{% we make the whole thing an ordinary symbol
		\left.\kern-\nulldelimiterspace % automatically resize the bar with \right
		#1 % the function
		\littletaller % pretend it's a little taller at normal size
		\right|_{#2} % this is the delimiter
}}

\newcommand{\littletaller}{\mathchoice{\vphantom{\big|}}{}{}{}}



%=========================================
%=========================================
%---For editing!!!!------------------------
%\newcommand{\mylabel}[1]{[#1]}
\newcommand{\mylabel}[1]{}

%\newcommand{\err}{{\bf !!!} \marginpar{\textcolor{red}{!!!}}}
%\newcommand{\err}{}

%\newcommand{\mycomment}[2]{ $\bf{\copyright_{#1}}$ \marginpar{\textcolor{red}{Comment 1:} #2}}
\newcommand{\mycomment}[1]{\marginpar{\addtolength{\baselineskip}{-2pt}{\footnotesize {\color{blue}  #1}}}}
%\newcommand{\mycomment}[1]{ {\textcolor{red}{Comment: #1}}}
%========================================
%========================================

% 1. Coloque esta linha no seu preâmbulo (antes do \begin{document}) 
% ou logo antes de abrir a sua matriz:
\setcounter{MaxMatrixCols}{25}

\title{Quantum multiplication matrices}
\author{}
\date{}

\begin{document}

\maketitle
\tableofcontents
\section{Complete intersections in $Fl(1,2,5)$}

\begin{enumerate}
    \item [1)]The Algebra is $A_4$, of $Sl_5$. We then exclude the last two simple roots and $\Pi_M = \{\alpha_1 , \alpha_2\}$. Therefore there are two generators in $H_2$, $\sigma_{1}, \sigma_2$.
    \item[2)] The first Chern class in this basis is $(2,4)$, i.e. $c_1 (M) = 2 \sigma_1 +4\sigma_2$.
    \item[3)]The Hodge diamond is
    \[
        \begin{array}{c}
             1  \\
             2\\
             3\\
             4\\
             4\\
            3\\
            2\\
            1
        \end{array}
    \]
\end{enumerate}


\subsection{Line bundle $\mathcal{O}(0,1) \oplus \mathcal{O}(0,2) \oplus \mathcal{O}(1,0)$}

\begin{enumerate}
    \item [1)] This the GM-20 in the Fano 4-fold of $K3$ type paper
    \item[2)]Birational to a general Gushei-Mukai fourfold
    \item[3)] Rationality not known according to the paper
    \item[4)] Hodge diamond is 
    \[
        \begin{array}{ccc}
             &1&  \\
             &2&\\
           1&3+25&1\\
            &2&\\
            &1&
        \end{array}
    \]
    Where 3+25 means $3$ ambient classes and $25$ primitive. 
\item[5)] First Chern class given by the coordinates $(1,1)$ and Fano index $1$. 

    \item[6)]

    \[
    \resizebox{\textwidth}{!}{$
\begin{pmatrix}
0 & 12y(1)y(2) & 4y(2)^2+12y(1)y(2) & 96y(1)y(2)^2 & 84y(1)y(2)^2 & 156y(1)y(2)^2 & 96y_{1}^{2}y_{2}^{2} +416 y(1) y(2)^3& 48y_{1}^{2}y_{2}^{2} +128 y(1) y(2)^3& 0 \\
1 & -y(1) & 6y(2) & 4y(2)^2 & 0 & 8y(2)^2 & 72y(1)y(2)^2 & 36y(1)y(2)^2 & 160 y(1) y(2)^3\\
1 & y(1) & 0 & 12y(1)y(2) & 20y(1)y(2) & 20y(1)y(2) & 132y(1)y(2)^2 & 48y(1)y(2)^2 & 48y_{1}^{2}y_{2}^{2} +128 y(1) y(2)^3\\
0 & 1 & 2 & 0 & 0 & 4y(2) & 8y(1)y(2) & 4y(1)y(2) & 36y(1)y(2)^2 \\
0 & 2 & 1 & 6y(2) & -y(1) & 10y(2) & 4y(2)^2 & 0 & 36y(1)y(2)^2 \\
0 & 0 & 1 & 0 & y(1) & -2y(2) & 16y(1)y(2) & 8y(1)y(2) & 24y(1)y(2)^2 \\
0 & 0 & 0 & 3 & 1 & 5 & y(1) & y(1) & 12y(1)y(2) \\
0 & 0 & 0 & -2 & 2 & -4 & 6y(2)-2y(1) & -2y(1) & 4y(2)^2-12y(1)y(2) \\
0 & 0 & 0 & 0 & 0 & 0 & 3 & 1 & 0
\end{pmatrix}
$}
\]

Let $y_1 = y_2 =1$
$$ 
\begin{pmatrix}
0 & 12 & 16 & 96 & 84 & 156 & 512 & 176 & 0 \\
1 & -1 & 6 & 4 & 0 & 8 & 72 & 36 & 160 \\
1 & 1 & 0 & 12 & 20 & 20 & 132 & 48 & 176 \\
0 & 1 & 2 & 0 & 0 & 4 & 8 & 4 & 36 \\
0 & 2 & 1 & 6 & -1 & 10 & 4 & 0 & 36 \\
0 & 0 & 1 & 0 & 1 & -2 & 16 & 8 & 24 \\
0 & 0 & 0 & 3 & 1 & 5 & 1 & 1 & 12 \\
0 & 0 & 0 & -2 & 2 & -4 & 4 & -2 & -8 \\
0 & 0 & 0 & 0 & 0 & 0 & 3 & 1 & 0
\end{pmatrix}
$$
\end{enumerate}

\section{Complete intersection in $\mathbb{P}^4$}

\subsection{Quadric 3-fold}

\begin{enumerate}
   
\item[1)] The hodge diamond is
\[
    \begin{array}{ccc}
         &1&  \\
          &1&  \\
           &1&  \\
            &1&  \\
    \end{array}
\]
\item[2)] The quantum multiplication table is
\[
    \begin{pmatrix}
        0& 0 &6y_1 &0\\
        3&0 &0 &6y_1\\
        0& 3 &0&0\\
        0&0&3&0
    \end{pmatrix}
\]
The eigenvalues are $\{0,3\sqrt{2y_1},3\sqrt{2y_1}\xi^{1/3},3\sqrt{2y_1}\xi^{2/3}$
\end{enumerate}


\subsection{Cubic 3-fold}
\begin{enumerate}
   

\item[1)] The hodge diamond is
\[
    \begin{array}{ccc}
         &1&  \\
          &1&  \\
          5&&5\\
           &1&  \\
            &1&  \\
    \end{array}
\]
\item[2)] The quantum multiplication table is
\[
    \begin{pmatrix}
        0& 12y_1 &0 &72y_{1}^{2}\\
        2&0 &30y_1 &0\\
        0& 2 &0&12y_1\\
        0&0&2&0
    \end{pmatrix}
\]
The eigenvalues are $\{0,0,6\sqrt{3y_1},-6\sqrt{3y_1}\}$
\end{enumerate}
\subsection{Quartic 3-fold}
\begin{enumerate}
 
\item[1)] The Hodge diamond is 
\[
    \begin{array}{ccc}
         &1&  \\
         &1&\\
         30&&30\\
         &1&\\
         &1&
    \end{array}
\]

\item[2)] The quantum multiplication matrix is
\[
    \begin{pmatrix}
        0& 3888y_{1}^{2} &504576y_{1}^{3} &18323712y_{1}^{4}\\
        1&80y_1 &13600y_{1}^{2} &504576y_{1}^{3}\\
        0& 1 &80y_1&3888y_{1}^{2}\\
        0&0&1&0
    \end{pmatrix}
 \]
 The eigenvalues are $-24y_1,-24y_1,-24y_1,232y_1$
 \item[3)] Irational. 
\end{enumerate}

\section{Complete intersection in Gr(2,4)}

\begin{enumerate}
    \item [1)]The Hodge diamond is
    \[
            \begin{array}{c}
             1  \\
             1\\
           2\\
            1\\
            1
        \end{array}
    \]

    \item[2)] The quantum multiplication table is 
    \[
    \begin{pmatrix}
        0 & 0 & 0 & 0 & 4y_2 & 0 \\
        4 & 0 & 0 & 0 & 0 & 4y_2 \\
        0 & 4 & 0 & 0 & 0 & 0 \\
        0 & 4 & 0 & 0 & 0 & 0 \\
        0 & 0 & 4 & 4 & 0 & 0 \\
        0 & 0 & 0 & 0 & 4 & 0
    \end{pmatrix}
    \]
There are $4$ roots of the unit and $0$ with multiplicity $2$.

\item[3)]There is a $\mathbb{Z}_2$ action that swaps the $2$ hodge classes of $H^4$ and leave $5$ invariants Hodge classes and the atomic decomposition is the one given by the small quantum multiplication. 
    
\end{enumerate}



\subsection{Complete intersection by the Line bundle $\mathcal{O}(1)$}

\begin{enumerate}
    \item [1)] This is the quadric $3$-fold because $Gr(2,4)$ is the quardic $4$-fold
\end{enumerate}

\subsection{Complete intersection by the Line bundle $\mathcal{O}(2)$}

\begin{enumerate}
    \item [1)] This is the del Pezzo $3$-4fold of degree $4$ and has Fano index 2. It is an intersection of two quadrics and is rational.

    \item[2)] The Hodge diamond is
    \[
    \begin{array}{ccc}
         &1&  \\
         &1&\\
         2&&2\\
         &1&\\
         &1&
    \end{array}
    \]
    \item[3)] The quantum multiplication matrix is
    \[
    \begin{pmatrix}
        0 & 8y_2 & 0 & 16y_2^2 \\
        2 & 0 & 8y_2 & 0 \\
        0 & 4 & 0 & 8y_2 \\
        0 & 0 & 2 & 0
    \end{pmatrix}
    \]
The eigenvalues are $\{0,0,8\sqrt{y_2},-8\sqrt{y_2}\}$
\end{enumerate}

\subsection{Complete intersection by the Line bundle $\mathcal{O}(3)$}
\begin{enumerate}
   

\item[1)] This is the complete intersection of a quadric and a cubic. Its suposed to be non-rational. It has Fano index 1
\item[2)] The Hodge diamond is 
\[
    \begin{array}{ccc}
         &1&  \\
         &1&\\
         20&&20\\
         &1&\\
         &1&
    \end{array}
\]
\item[3)] The quantum multiplication table is
\[
\begin{pmatrix}
0 & 792y_2^{2} & 21816y_2^{3} & 396576y_2^{4} \\
1 & 30y_2 & 1170y_2^{2} & 21816y_2^{3} \\
0 & 2 & 30y_2 & 792y_2^{2} \\
0 & 0 & 1 & 0
\end{pmatrix}
\]
The eigenvalues are $\{-12y_2 , -12y_2 , -12y_2, 96y_2\}$.
\end{enumerate}
\section{Complete intersection in $\mathbb{P}^5$}


\subsection{Quadric 4-fold}

\begin{enumerate}
    \item [1)] The quadric 4-fold is the $Gr(2,4)$

\item[2)] The Hodge diamond is
\[
    \begin{array}{c}
         1  \\
          1\\
          2\\
          1\\
          1
    \end{array}
\]
    
    \item[3)] The quantum multiplication matrix is
    \[
    \begin{pmatrix}
        0 & 0 & 0 & 8y_1 & 0 \\
        4 & 0 & 0 & 0 & 8y_1 \\
        0 & 4 & 0 & 0 & 0 \\
        0 & 0 & 4 & 0 & 0 \\
        0 & 0 & 0 & 4 & 0
    \end{pmatrix}
    \]
The eigenvalues are $\{0,\pm 4 \sqrt{2}y_{1}^{1/4} ,\pm 4i \sqrt{2}y_{1}^{1/4} \}$
\end{enumerate}


\subsection{Cubic 4-fold}

\begin{enumerate}
    \item [1)] The Hodge diamond is
    \[
        \begin{array}{ccccc}
             &1&  \\
             &1&\\
             1&21&1\\
             &1&\\
             &1&
        \end{array}
    \]

    \item[2)] This is non-rational

    \item[3)] The quantum multiplication matrix is 



\[
\begin{pmatrix}
0 & 0 & 18 y_1 & 0 & 0 \\
3 & 0 & 0 & 45 y_1 & 0 \\
0 & 3 & 0 & 0 & 18 y_1 \\
0 & 0 & 3 & 0 & 0 \\
0 & 0 & 0 & 3 & 0
\end{pmatrix}
\]
Eigenvalues $0,0,9\xi y_1^{1/3}, 9 \xi^{1/3} y_1^{1/3},
 9\xi^{2/3} y_1^{1/3}$
 \end{enumerate}

\subsection{Quartic 4-fold}
\begin{enumerate}

\item[1)]The Hodge diamond is
\[
    \begin{array}{ccc}
         &1&  \\
         &1&\\
         21 &1+141 &21\\
         &1&\\
         &1&
    \end{array}
\]
where $1+141$ means $1$ class is in the ambient cohomology and $141$ in the primitive.
\item[2)] The quantum multiplication matrix is 
\[
\begin{pmatrix}
0 & 48 y_1 & 0 & 5568 y_1^2 & 0 \\
2 & 0 & 208 y_1 & 0 & 5568 y_1^2 \\
0 & 2 & 0 & 208 y_1 & 0 \\
0 & 0 & 2 & 0 & 48 y_1 \\
0 & 0 & 0 & 2 & 0
\end{pmatrix}
\]

Eigenvalues are $0,0,0,-32\sqrt{y_1},32\sqrt{y_1}$
\end{enumerate}



\subsection{Intersection of two quadrics $\mathcal{O}(2) \oplus \mathcal{O}(2)$}

\begin{itemize}
\item[1)] The Hodge diamond is
\[
    \begin{array}{ccc}
         &1  \\
         &1\\
         2&&2\\
         &1\\
         &1
    \end{array}
\]


\item[2)] The quantum multiplication table is

\[
\begin{pmatrix}
0 & 8y_1 & 0 & 32y_1^2 \\
2 & 0 & 16y_1 & 0 \\
0 & 2 & 0 & 8y_1 \\
0 & 0 & 2 & 0
\end{pmatrix}
\]

The eigenvalues are $\{0,0,8\sqrt{y(1)},-8\sqrt{y_1}\}$


\end{itemize}

\section{$Gr(2,6)$}

\subsection{$\mathcal{O}(1)\oplus \mathcal{O}(1) \oplus \mathcal{O}(2)$}
\[
\begin{array}{ccccccccccc}
&&&&&1&&&&&\\
&&&&0&&0&&&&\\
&&&0&&1&&0&&&\\
&&0&&0&&0&&0&&\\
&0&&0&&2&&0&&0&\\
0&&1&&67&&67&&1&&0\\
&0&&0&&2&&0&&0&\\
&&0&&0&&0&&0&&\\
&&&0&&1&&0&&&\\
&&&&0&&0&&&&\\
&&&&&1&&&&&
\end{array}
\]
\[
\begin{pmatrix}
0 & 16y_{2} & 0 & 0 & 336y_{2}^{2} & 288y_{2}^{2} & 0 & 768y_{2}^{3} \\
2 & 0 & 16y_{2} & 36y_{2} & 0 & 0 & 480y_{2}^{2} & 0 \\
0 & 2 & 0 & 0 & 12y_{2} & 12y_{2} & 0 & 48y_{2}^{2} \\
0 & 2 & 0 & 0 & 32y_{2} & 28y_{2} & 0 & 80y_{2}^{2} \\
0 & 0 & 2 & 2 & 0 & 0 & 12y_{2} & 0 \\
0 & 0 & 0 & 2 & 0 & 0 & 24y_{2} & 0 \\
0 & 0 & 0 & 0 & \frac{10}{3} & \frac{8}{3} & 0 & \frac{16y_{2}}{3} \\
0 & 0 & 0 & 0 & 0 & 0 & 6 & 0
\end{pmatrix}
\]
\[
\left(
0,\,
0,\,
0,\,
0,\,
-4i\sqrt{y_{2}},\,
4i\sqrt{y_{2}},\,
-12\sqrt{3}\sqrt{y_{2}},\,
12\sqrt{3}\sqrt{y_{2}}
\right)
\]

\section{Complete intersection in $Gr(2,5)$}



\subsection{Line bundle $\mathcal{O}(1)$, A linear section}
\begin{enumerate}

\item[1)]The Hodge diamond is 
\[
    \begin{array}{c}
         1  \\
          1\\
          2\\
          2\\
          1\\
          1
    \end{array}
\]

\item[2)]  The quantum multiplication table is
\[
M=
\begin{pmatrix}
0&0&0&0&4y_1&4y_1&0&0\\
4&0&0&0&0&0&4y_1&0\\
0&4&0&0&0&0&0&0\\
0&4&0&0&0&0&0&4y_1\\
0&0&4&4&0&0&0&0\\
0&0&0&4&0&0&0&0\\
0&0&0&0&8&4&0&0\\
0&0&0&0&0&0&4&0
\end{pmatrix}.
\]
The eigenvalues are
\[
\lambda \in 
\left\{
\zeta_4^k\sqrt[4]{128y_1\bigl(11+5\sqrt5\bigr)},
\;
\zeta_4^k\sqrt[4]{128y_1\bigl(11-5\sqrt5\bigr)}
\;\middle|\;
k=0,1,2,3
\right\},
\qquad \zeta_4=i.
\]
\item[3)] This is a del-Pezzo $5$-fold $V_5$. It is rational. 
\end{enumerate}
\subsection{Line bundle $\mathcal{O}(2)$ Gushel-Mukai 5-fold}

\begin{enumerate}
    
\item[1)] The Hodge diamond is
\[
\begin{array}{cccccccccccccc}
&& && && 1 && && && \\
&& && & 0 && 0 & && && \\
&& && 0 && 1 && 0 && && \\
&& & 0 && 0 && 0 && 0 && & \\
&& 0 && 0 && 2 && 0 && 0 && \\
& 0 && 0 && 10 && 10 && 0 && 0  \\
&& 0 && 0 && 2 && 0 && 0 && \\
&& & 0 && 0 && 0 && 0 && & \\
&& && 0 && 1 && 0 && && \\
&& && & 0 && 0 & && && \\
&& && && 1 && && && \\
\end{array}
\]
\item[2)]
\[
M=
\begin{pmatrix}
0&0&6y_1&12y_1&0&0&0&24y_1^2\\
3&0&0&0&18y_1&12y_1&0&0\\
0&3&0&0&0&0&6y_1&0\\
0&3&0&0&0&0&12y_1&0\\
0&0&3&3&0&0&0&6y_1\\
0&0&0&3&0&0&0&6y_1\\
0&0&0&0&6&3&0&0\\
0&0&0&0&0&0&3&0
\end{pmatrix}.
\]
The eigenvalues are
\[
\left\{
0,0,
\omega^k\sqrt[3]{54y_1(11+5\sqrt5)},
\omega^k\sqrt[3]{54y_1(11-5\sqrt5)}
\;:\; k=0,1,2
\right\}.
\]
\item[3)] A quadric section in $Gr(2,5)$, It is the Gushel-Mukai 5-fold. It is rational. 
\end{enumerate}

\subsection{Line bundle $\mathcal{O}(3)$}
\begin{enumerate}
    \item [1)] The Hodge diamond is
    \[
\begin{array}{ccccccccccc}
 &&&&&1&&&&&\\
 &&&&0&&0&&&&\\
 &&&0&&1&&0&&&\\
 &&0&&0&&0&&0&&\\
 &0&&0&&2&&0&&0&\\
0&&10&&300&&300&&10&&0\\
 &0&&0&&2&&0&&0&\\
 &&0&&0&&0&&0&&\\
 &&&0&&1&&0&&&\\
 &&&&0&&0&&&&\\
 &&&&&1&&&&&
\end{array}
\]

\item[2)] The quantum multiplication table is 
\[
\begin{pmatrix}
0 & 36y(2) & 0 & 0 & 2232y(2)^2 & 1224y(2)^2 & 0 & 26568y(2)^3 \\
2 & 0 & 54y(2) & 96y(2) & 0 & 0 & 2988y(2)^2 & 0 \\
0 & 2 & 0 & 0 & 72y(2) & 42y(2) & 0 & 1008y(2)^2 \\
0 & 2 & 0 & 0 & 96y(2) & 54y(2) & 0 & 1224y(2)^2 \\
0 & 0 & 2 & 2 & 0 & 0 & 54y(2) & 0 \\
0 & 0 & 0 & 2 & 0 & 0 & 42y(2) & 0 \\
0 & 0 & 0 & 0 & 4 & 2 & 0 & 36y(2) \\
0 & 0 & 0 & 0 & 0 & 0 & 2 & 0
\end{pmatrix}
\]
The characteristic polynomial is
\[
\chi(\lambda)=\lambda^8-1188\,y(2)\,\lambda^6-11664\,y(2)^2\,\lambda^4
=\lambda^4\Bigl(\lambda^4-1188\,y(2)\,\lambda^2-11664\,y(2)^2\Bigr).
\]

Hence the eigenvalues are
\[
0 \quad (\text{with multiplicity }4),
\]
and
\[
\lambda=\pm \sqrt{54\,y(2)\,\bigl(11+5\sqrt5\bigr)},\qquad
\lambda=\pm \sqrt{54\,y(2)\,\bigl(11-5\sqrt5\bigr)}.
\]
\end{enumerate}
\subsection{Line bundle $\mathcal{O}(1) \oplus \mathcal{O}(1)$}

\begin{enumerate}
    \item [1)]
    \[
    \begin{array}{c}
         1  \\
          1\\
          2\\
          1\\
          1
    \end{array}
    \]

\item [2)]  The small quantum multiplication table is
\[
M=
\begin{pmatrix}
0&0&3y_1&6y_1&0&0\\
3&0&0&0&6y_1&0\\
0&3&0&0&0&0\\
0&3&0&0&0&3y_1\\
0&0&3&\frac92&0&0\\
0&0&0&0&6&0
\end{pmatrix}.
\]
The eigenvalues are
\[
\lambda=\omega^k\sqrt[3]{\frac{y_1}{2}\left(297+135\sqrt5\right)}
\quad\text{and}\quad
\lambda=\omega^k\sqrt[3]{\frac{y_1}{2}\left(297-135\sqrt5\right)},
\qquad k=0,1,2,
\]
where
\[
\omega=e^{2\pi i/3}.
\]

\item[3)] This is a del-Pezzo $4$-fold. It is rational
\end{enumerate}

\subsection{Line bundle $\mathcal{O}(1) \oplus \mathcal{O}(2)$ Gushel-Mukai 4-fold}

\begin{enumerate}
\item[1)] This is the Gushel-Mukai 4-fold. The very general is known to be non-rational.
    \item [2)]The Hodge diamond is
    \[
        \begin{array}{ccc}
             &1&  \\
             &1&\\
             1&22&1\\
             &1&\\
             &1&
        \end{array}
    \]

    \item[3)] The matrix of the quantum multiplication is

    \[
M=
\begin{pmatrix}
0&12y_1&0&0&96y_1^2&0\\
2&0&12y_1&20y_1&0&48y_1^2\\
0&2&0&0&8y_1&0\\
0&2&0&0&16y_1&0\\
0&0&2&3&0&6y_1\\
0&0&0&0&4&0
\end{pmatrix}.
\]
\begin{align}
&\lambda = 0, \ 0, \ \pm 2\sqrt{2y_1(11 + 5\sqrt{5})}, \ \pm 2\sqrt{2y_1(11 - 5\sqrt{5})}
\end{align}
% Compact simplified form

\end{enumerate}

\subsection{Line bundle $\mathcal{O}(1) \oplus \mathcal{O}(1) \oplus \mathcal{O}(1)$}

\begin{enumerate}
    \item [1)]The Hodge diamond is
    \[
        \begin{array}{c}
             1  \\
             1\\
             1\\
             1
        \end{array}
    \]

\[
M=
\begin{pmatrix}
0&6y_1&0&8y_1^2\\
2&0&4y_1&0\\
0&5&0&6y_1\\
0&0&2&0
\end{pmatrix}.
\]
The eigenvalues are
\[
    \lambda = \pm \sqrt{2y_1(11 + 5\sqrt{5})}, \ \pm \sqrt{2y_1(11 - 5\sqrt{5})}
\]
\item[3)] This is the del-Pezzo 3-fold $V_5$. It is rational 
\end{enumerate}

\section{Complete intersection in $\mathbb{P}^6$}

\subsection{Quadric 5-fold}

\begin{enumerate}
    \item [1)]The hodge diamond is
    \[
        \begin{array}{c}
            1   \\
             1\\
             1\\
             1\\
             1\\
             1
        \end{array}
    \]

    \item[2)] The quantum multiplication table is
    \[
A=
\begin{pmatrix}
0 & 0 & 0 & 0 & 10y_1 & 0 \\
5 & 0 & 0 & 0 & 0 & 10y_1 \\
0 & 5 & 0 & 0 & 0 & 0 \\
0 & 0 & 5 & 0 & 0 & 0 \\
0 & 0 & 0 & 5 & 0 & 0 \\
0 & 0 & 0 & 0 & 5 & 0
\end{pmatrix}.
\]
And the eigenvalues
\[
\lambda_0=0,
\qquad
\lambda_k=\sqrt[5]{12500\,y_1}\,\zeta_5^k
\quad (k=0,1,2,3,4),
\]
\end{enumerate}

\subsection{Cubic 5-fold}

\begin{enumerate}
  \item [1)]The hodge diamond is
    \[
        \begin{array}{ccc}
            &1&   \\
             &1&\\
             &1&\\
             21&&21\\
             &1&\\
             &1&\\
             &1&
        \end{array}
    \]
    \item[] The quantum multiplication table is

    \[
A=
\begin{pmatrix}
0 & 0 & 0 & 24y_1 & 0 & 0 \\
4 & 0 & 0 & 0 & 60y_1 & 0 \\
0 & 4 & 0 & 0 & 0 & 24y_1 \\
0 & 0 & 4 & 0 & 0 & 0 \\
0 & 0 & 0 & 4 & 0 & 0 \\
0 & 0 & 0 & 0 & 4 & 0
\end{pmatrix}.
\]

And the eigenvalues are
\[
\lambda_1=\lambda_2=0,
\]
\[
\lambda_k
=
\sqrt[4]{6912\,y_1}\,i^{\,k},
\qquad k=0,1,2,3.
\]
\end{enumerate}

\subsection{Quartic 5-fold}
The Hodge diamond is 
\[
    \begin{array}{ccccc}
         &&1&&  \\
         &&1&&\\
         &&1&&\\
         7&266&&266&7\\
            &&1&&  \\
         &&1&&\\
         &&1&&
    \end{array}
\]
\[
\begin{pmatrix}
0 & 0 & 72\, y_1 & 0 & 0 & 1728\, y_1^2 \\
3 & 0 & 0 & 312\, y_1 & 0 & 0 \\
0 & 3 & 0 & 0 & 312\, y_1 & 0 \\
0 & 0 & 3 & 0 & 0 & 72\, y_1 \\
0 & 0 & 0 & 3 & 0 & 0 \\
0 & 0 & 0 & 0 & 3 & 0
\end{pmatrix}
\]
Eigenvalues are $\{0,0,0,12(-2)^{2/3}\, y_1^{1/3}, \quad
12 \cdot 2^{2/3}\, y_1^{1/3}, \quad
-12(-1)^{1/3} 2^{2/3}\, y_1^{1/3}\}$\\
It is irational
\subsection{Quintic 5-fold}

The Hodge diamond is 
\[
    \begin{array}{ccccc}
         &&1&&  \\
         &&1&&\\
         &&1&&\\
         84&1554&&1554&84\\
            &&1&&  \\
         &&1&&\\
         &&1&&
    \end{array}
\]
\[
\begin{pmatrix}
0 & 240 y_1 & 0 & 422400 y_1^2 & 0 & 62640000 y_1^3 \\
2 & 0 & 1540 y_1 & 0 & 1385000 y_1^2 & 0 \\
0 & 2 & 0 & 2690 y_1 & 0 & 422400 y_1^2 \\
0 & 0 & 2 & 0 & 1540 y_1 & 0 \\
0 & 0 & 0 & 2 & 0 & 240 y_1 \\
0 & 0 & 0 & 0 & 2 & 0
\end{pmatrix}
\]
Eigenvalues $0,0,0,0,-50\sqrt{5y_1} ,50\sqrt{5y_1}$\\
It is irrational.
\subsection{Intersection of quadric and cubic}

\begin{enumerate}
    \item [1)] The Hodge diamond is
    \[
\begin{array}{ccc}
     &1&  \\
     &1&\\
     8&70&8\\
          &1&  \\
     &1&
\end{array}
    \]
    \item[2)]  The quantum multiplication table is
    \[
\begin{pmatrix}
0 & 24y(1) & 0 & 1152y(1)^2 & 0 \\
2 & 0 & 84y(1) & 0 & 1152y(1)^2 \\
0 & 2 & 0 & 84y(1) & 0 \\
0 & 0 & 2 & 0 & 24y(1) \\
0 & 0 & 0 & 2 & 0
\end{pmatrix}
\]
The eigenvalues are $\{0,0,0,12\sqrt{3}y_1 , -12\sqrt{3}y_1\}$
\end{enumerate}


\section{Complete intersection in $Fl(1,4,5)$}

\begin{enumerate}
\item[1)] This should be a complete intersection of $\mathbb{P}^4 \times \mathbb{P}^4$ by the line bundle $\mathcal{O}(1,1)$.
    \item [2)] The Hodge diamond is 
    \[
\begin{array}{c}
1\\
2\\
3\\
4\\
4\\
3\\
2\\
1
\end{array}
\]
    \item[3)] The quantum multiplication matrix is

% \[
% {\begin{pmatrix}
% 0&0&0&0&0&0&0&0&4y_2&4y_1&0&0&0&0&0&0&0&0&0&8y_1 y_2\\
% 4&0&0&0&0&0&0&0&0&0&4y_2&0&0&0&0&0&0&0&0&0\\
% 4&0&0&0&0&0&0&0&0&0&0&0&0&4y_1&0&0&0&0&0&0\\
% 0&4&4&0&0&0&0&0&0&0&0&0&0&0&0&0&0&0&0&0\\
% 0&4&0&0&0&0&0&0&0&0&0&0&0&0&4y_2&0&0&0&0&0\\
% 0&0&4&0&0&0&0&0&0&0&0&0&0&0&0&0&4y_1&0&0&0\\
% 0&0&0&4&0&4&0&0&0&0&0&0&0&0&0&0&0&0&0&0\\
% 0&0&0&4&4&0&0&0&0&0&0&0&0&0&0&0&0&0&0&0\\
% 0&0&0&0&0&4&0&0&0&0&0&0&0&0&0&0&0&0&4y_1&0\\
% 0&0&0&0&4&0&0&0&0&0&0&0&0&0&0&0&0&4y_2&0&0\\
% 0&0&0&0&0&0&4&0&8&0&0&0&0&0&0&0&0&0&0&4y_1\\
% 0&0&0&0&0&0&8&4&4&0&0&0&0&0&0&0&0&0&0&0\\
% 0&0&0&0&0&0&4&8&0&4&0&0&0&0&0&0&0&0&0&0\\
% 0&0&0&0&0&0&0&4&0&8&0&0&0&0&0&0&0&0&0&4y_2\\
% 0&0&0&0&0&0&0&0&0&0&4&4&0&0&0&0&0&0&0&0\\
% 0&0&0&0&0&0&0&0&0&0&0&4&4&0&0&0&0&0&0&0\\
% 0&0&0&0&0&0&0&0&0&0&0&0&4&4&0&0&0&0&0&0\\
% 0&0&0&0&0&0&0&0&0&0&0&0&0&0&4&4&0&0&0&0\\
% 0&0&0&0&0&0&0&0&0&0&0&0&0&0&0&4&4&0&0&0\\
% 0&0&0&0&0&0&0&0&0&0&0&0&0&0&0&0&0&4&4&0
% \end{pmatrix}}
% \]


% 2. Agora o LaTeX vai aceitar as 20 colunas sem quebrar nada:
\[
\begin{pmatrix}
0&0&0&0&0&0&0&0&4y_2&4y_1&0&0&0&0&0&0&0&0&0&8y_1 y_2\\
4&0&0&0&0&0&0&0&0&0&4y_2&0&0&0&0&0&0&0&0&0\\
4&0&0&0&0&0&0&0&0&0&0&0&0&4y_1&0&0&0&0&0&0\\
0&4&4&0&0&0&0&0&0&0&0&0&0&0&0&0&0&0&0&0\\
0&4&0&0&0&0&0&0&0&0&0&0&0&0&4y_2&0&0&0&0&0\\
0&0&4&0&0&0&0&0&0&0&0&0&0&0&0&0&4y_1&0&0&0\\
0&0&0&4&0&4&0&0&0&0&0&0&0&0&0&0&0&0&0&0\\
0&0&0&4&4&0&0&0&0&0&0&0&0&0&0&0&0&0&0&0\\
0&0&0&0&0&4&0&0&0&0&0&0&0&0&0&0&0&0&4y_1&0\\
0&0&0&0&4&0&0&0&0&0&0&0&0&0&0&0&0&4y_2&0&0\\
0&0&0&0&0&0&4&0&8&0&0&0&0&0&0&0&0&0&0&4y_1\\
0&0&0&0&0&0&8&4&4&0&0&0&0&0&0&0&0&0&0&0\\
0&0&0&0&0&0&4&8&0&4&0&0&0&0&0&0&0&0&0&0\\
0&0&0&0&0&0&0&4&0&8&0&0&0&0&0&0&0&0&0&4y_2\\
0&0&0&0&0&0&0&0&0&0&4&4&0&0&0&0&0&0&0&0\\
0&0&0&0&0&0&0&0&0&0&0&4&4&0&0&0&0&0&0&0\\
0&0&0&0&0&0&0&0&0&0&0&0&4&4&0&0&0&0&0&0\\
0&0&0&0&0&0&0&0&0&0&0&0&0&0&4&4&0&0&0&0\\
0&0&0&0&0&0&0&0&0&0&0&0&0&0&0&4&4&0&0&0\\
0&0&0&0&0&0&0&0&0&0&0&0&0&0&0&0&0&4&4&0
\end{pmatrix}
\]

The eigenvalues are
\[
a:=y(1),\qquad b:=y(4),\qquad \omega:=i.
\]

\noindent\textbf{Case \(y_1\neq y_2\).}
Let \(z=\lambda^4\). Then the characteristic polynomial is
\[
Q(z)=z^5-1280(y_1+y_2)z^4+\bigl(655360(y_1^2+y_2^2)-39649280ab\bigr)z^3
\]
\[
\hspace{2.2cm}
-\bigl(167772160(y_1^3+y_2^3)+31960596480y_1 y_2(y_1+y_2)\bigr)z^2
\]
\[
\hspace{2.2cm}
+\bigl(21474836480(y_1^4+y_2^4)-2598455214080y_1 y_2(y_1^2+y_2^2)+8181912698880y_1^2y_2^2\bigr)z
\]
\[
\hspace{2.2cm}
-1099511627776(y_1^5+y_2^5)-5497558138880(y_1^4 y_2+y_1 y_2^4)-10995116277760(y_1^3 y_2^2 + y_1^2 y_2^3).
\]
If \(z_1,\dots,z_5\) are the roots of \(Q(z)\), then the eigenvalues are
\[
\lambda=\omega^k z_j^{1/4},\qquad j=1,\dots,5,\quad k=0,1,2,3.
\]

\medskip

\noindent\textbf{Case \(y_1=y_2=t\).}
Then
\[
\chi_M(\lambda)
=
(\lambda^4-8192t)\,\bigl(\lambda^8+2816t\,\lambda^4-65536t^2\bigr)^2.
\]
Equivalently,
\[
\lambda^4=8192t,
\]
or
\[
\lambda^4=128t(5\sqrt5-11),
\qquad
\lambda^4=-128t(11+5\sqrt5).
\]
Hence the eigenvalues are
\[
\lambda=\omega^k(8192t)^{1/4},\qquad k=0,1,2,3,
\]
\[
\lambda=\omega^k\bigl(128t(5\sqrt5-11)\bigr)^{1/4},\qquad k=0,1,2,3,
\]
\[
\lambda=\omega^k\bigl(-128t(11+5\sqrt5)\bigr)^{1/4},\qquad k=0,1,2,3.
\]


\item[3)] We have 20 different eigenvalues if $y_1 \neq y_2$ and 12 if $y_1=y_2$. Again this happens because we have a $Z_2$ action in $Fl(1,4,5)$ that has $12$ invariant Hodge classes. 
\end{enumerate}


\subsection{Line bundle $\mathcal{O}(1,1) \oplus \mathcal{O}(1,1)$}
\begin{enumerate}
    \item [1)] The Hodge diamond is
    \[
        \begin{array}{ccc}
             &1&  \\
             &2&\\
             &3&\\
             6&&6\\
             &3&\\
             &2&\\
             &1&
        \end{array}
    \]
\item[2)] The quantum multiplication atrix is
\[
a:=y(1),\qquad b:=y(4),
\]
\[
\begin{pmatrix}
0&2y_1&2y_2&0&0&0&28y_1 y_2&28y_1 y_2&16y_1 y_2&0&0&6y_2 y_1^2+6y_2^2 y_1\\
2&0&0&2y_2&0&6y_2&0&0&0&4y_1 y_2&16y_1 y_2&0\\
2&0&0&2y_1&6y_1&0&0&0&0&8y_1 y_2&20y_1 y_2&0\\
0&2&2&0&0&0&0&0&0&0&0&4y_1 y_2\\
0&2&0&0&0&0&6 y_2&2y_2&6y_2&0&0&4y_1 y_2\\
0&0&2&0&0&0&2 y_1&6 y_1&0&0&0&4y_1 y_2\\
0&0&0&2&-3&2&0&0&0&-3y_2&-6y_2&0\\
0&0&0&2&5&0&0&0&0&3y_2&6y_2&0\\
0&0&0&0&2&2&0&0&0&2y_2&2y_1+4y_2&0\\
0&0&0&0&0&0&-4&-16&4&0&0&2y_1-4y_2\\
0&0&0&0&0&0&8&12&2&0&0&2y_2\\
0&0&0&0&0&0&0&0&0&2&6&0
\end{pmatrix}.
\]

The characteristic polynomial is
\[
\chi_M(\lambda)=\lambda^2\,Q(\lambda^2),
\]
where
\[
\begin{aligned}
Q(z)=\;&z^5-20sz^4+\bigl(160s^2-10000p\bigr)z^3
-\bigl(640s^3+120000sp\bigr)z^2\\
&+\bigl(1280s^4-160000s^2p+800000p^2\bigr)z-1024s^5.
\end{aligned}
\]

Hence the eigenvalues are:
\[
\lambda=0 \quad (\text{with multiplicity }2),
\]
and, if \(z_1,\dots,z_5\) are the roots of \(Q(z)\),
\[
\lambda=\pm \sqrt{z_1},\ \pm \sqrt{z_2},\ \pm \sqrt{z_3},\ \pm \sqrt{z_4},\ \pm \sqrt{z_5}.
\]
\[
a:=y(1),\qquad b:=y(4),\qquad s:=a+b,\qquad p:=ab,\qquad z:=\lambda^2.
\]
\[
a:=y(1),\qquad b:=y(4),\qquad s:=a+b,\qquad p:=ab,\qquad z:=\lambda^2.
\]
If $a=b$ we have
\[
\chi_M(\lambda)=\lambda^2\bigl(\lambda^2-128t\bigr)\bigl(\lambda^4+44t\,\lambda^2-16t^2\bigr)^2,
\]


\end{enumerate}


\section{Complete intersection in $B_2$}

\subsection{Line bundle $\mathcal{O}(1,1)$}


\begin{enumerate}
    \item ]1) The Hodge diamond is
\[
        \begin{array}{ccc}
             &1&  \\
             &2&\\
             1&&1\\
             &2&\\
             &1&
        \end{array}
  \]
\item[2)]  \[
  \resizebox{\textwidth}{!}{$
\begin{pmatrix}
0 & 10y(1)y(2) & 10y(1)y(2) & 15y(2)y(1)^2+9y(2)^2y(1) & 15y(2)y(1)^2+3y(2)^2y(1) & 4y(2)y(1)^3+20y(2)^2y(1)^2 \\
1 & -y(1) & 2y(2) & 8y(1)y(2) & 6y(1)y(2) & 3y(2)y(1)^2+3y(2)^2y(1) \\
1 & 3y(1) & -y(2) & 6y(1)y(2) & 2y(1)y(2) & 6y(2)y(1)^2 \\
0 & 1 & 2 & -y(2) & 2y(1) & 2y(1)y(2) \\
0 & 3 & 1 & 3y(2) & -y(1) & 4y(1)y(2) \\
0 & 0 & 0 & 4 & 3 & 0
\end{pmatrix}
$}
\]

\[
\chi_M(\lambda)
=
-(\lambda+y(1)+y(2))^2
\Bigl(
-\lambda^4
+90\,\lambda^2 y(1)y(2)
+320\,\lambda\, y(1)^2y(2)
+108\,\lambda\, y(1)y(2)^2
+256\,y(1)^3y(2)
+135\,y(1)^2y(2)^2
\Bigr).
\]

Hence the eigenvalues are
\[
\lambda=-y(1)-y(2)
\quad\text{with multiplicity }2,
\]
and the remaining four eigenvalues are the roots of
\[
-\lambda^4
+90\,\lambda^2 y(1)y(2)
+320\,\lambda\, y(1)^2y(2)
+108\,\lambda\, y(1)y(2)^2
+256\,y(1)^3y(2)
+135\,y(1)^2y(2)^2
=0.
\]
\end{enumerate}


\section{$G_2/P_1$ the quadric 5-fold}


\begin{enumerate}
    \item [1)]The hodge diamond is
    \[
        \begin{array}{c}
            1   \\
             1\\
             1\\
             1\\
             1\\
             1
        \end{array}
    \]

    \item[2)] The quantum multiplication table is
    \[
A=
\begin{pmatrix}
0 & 0 & 0 & 0 & 5y_1 & 0 \\
5 & 0 & 0 & 0 & 0 & 5y_1 \\
0 & 5 & 0 & 0 & 0 & 0 \\
0 & 0 & 10 & 0 & 0 & 0 \\
0 & 0 & 0 & 5 & 0 & 0 \\
0 & 0 & 0 & 0 & 5 & 0
\end{pmatrix}.
\]
And the eigenvalues
\[
\lambda_0=0,
\qquad
\lambda_k=\sqrt[5]{12500\,y_1}\,\zeta_5^k
\quad (k=0,1,2,3,4),
\]
\end{enumerate}

\section{Complete intersection in $G_2/P_2$}
The hodge diamond is
    \[
        \begin{array}{c}
            1   \\
             1\\
             1\\
             1\\
             1\\
             1
        \end{array}
    \]

\[
A=
\begin{pmatrix}
0 & 0 & 3y(2) & 0 & 0 & 6y(2)^2\\
3 & 0 & 0 & 3y(2) & 0 & 0\\
0 & 9 & 0 & 0 & 3y(2) & 0\\
0 & 0 & 6 & 0 & 0 & 3y(2)\\
0 & 0 & 0 & 9 & 0 & 0\\
0 & 0 & 0 & 0 & 3 & 0
\end{pmatrix}.
\]

The characteristic polynomial is
\[
\chi_A(\lambda)=\lambda^6-486\,y(2)\lambda^3-19683\,y(2)^2
=
\bigl(\lambda^3-81(3+2\sqrt3)\,y(2)\bigr)
\bigl(\lambda^3-81(3-2\sqrt3)\,y(2)\bigr).
\]

Hence the eigenvalues are the six cube roots of the two numbers
\(81(3\pm 2\sqrt3)\,y(2)\). Equivalently,
\[
\lambda=\omega^k \sqrt[3]{81(3+2\sqrt3)\,y(2)}
\quad\text{and}\quad
\lambda=\omega^k \sqrt[3]{81(3-2\sqrt3)\,y(2)},
\qquad k=0,1,2,
\]
where \(\omega=e^{2\pi i/3}\).

If you want them written in the same Mathematica style as in your screenshot, you can write:
\[
\left\{
-3\sqrt[3]{-3}\sqrt[3]{(3-2\sqrt3)\,y(2)},\;
3\sqrt[3]{3}\sqrt[3]{(3-2\sqrt3)\,y(2)},\;
3(-1)^{2/3}\sqrt[3]{3}\sqrt[3]{(3-2\sqrt3)\,y(2)},
\right.
\]
\[
\left.
-3\sqrt[3]{-3}\sqrt[3]{(3+2\sqrt3)\,y(2)},\;
3\sqrt[3]{3}\sqrt[3]{(3+2\sqrt3)\,y(2)},\;
3(-1)^{2/3}\sqrt[3]{3}\sqrt[3]{(3+2\sqrt3)\,y(2)}
\right\}.
\]



\subsection{Line bundle $\mathcal{O}(1)$}

The hodge diamond is
    \[
        \begin{array}{c}
            1   \\
             1\\
             1\\
             1\\
             1
        \end{array}
    \]

\[
A=
\begin{pmatrix}
0 & 6y\left(2\right) & 0 & 12y\left(2\right)^2 & 0 \\
2 & 0 & 4y(2) & 0 & 4y(2)^2 \\
0 & 6 & 0 & 6y(2) & 0 \\
0 & 0 & 4 & 0 & 2y(2) \\
0 & 0 & 0 & 6 & 0
\end{pmatrix}.
\]

Its eigenvalues are
\[
\left\{
0,\,
-2\sqrt3\,\sqrt{\,3y(2)-2\sqrt3\,y(2)\,},
\,
2\sqrt3\,\sqrt{\,3y(2)-2\sqrt3\,y(2)\,},
\,
-2\sqrt3\,\sqrt{\,2\sqrt3\,y(2)+3y(2)\,},
\,
2\sqrt3\,\sqrt{\,2\sqrt3\,y(2)+3y(2)\,}
\right\}.
\]

Equivalently, the characteristic polynomial is
\[
\chi_A(\lambda)
=
\lambda^5-72\,y(2)\lambda^3-432\,y(2)^2\lambda
=
\lambda\bigl(\lambda^2-12(3-2\sqrt3)\,y(2)\bigr)
\bigl(\lambda^2-12(3+2\sqrt3)\,y(2)\bigr).
\]


\section{Complete intersection in $Fl(1,3,4)$}
This variety is $\mathbb{P}^3 \times \mathbb{P}^3$ cut out by the line bundle $\mathcal{O}(1,1)$.
\subsection{Line bundle $\mathcal{O}(1,1)$}

The hodge diamond is
\[
    \begin{array}{c}
         1  \\
          2\\
          6\\
          2\\
          1
    \end{array}
\]
where $3$ out of the $6$ are ambient
\[
  \resizebox{\textwidth}{!}{$
\begin{pmatrix}
0 & 2y(1) & 2y(3) & 0 & 0 & 0 & 4y(1)y(3) & 4y(1)y(3) & 0 \\
2 & 0 & 0 & 2y(3) & 0 & 4y(3) & 0 & 0 & 4y(1)y(3) \\
2 & 0 & 0 & 2y(1) & 4y(1) & 0 & 0 & 0 & 4y(1)y(3) \\
0 & 2 & 2 & 0 & 0 & 0 & 0 & 0 & 0 \\
0 & 2 & 0 & 0 & 0 & 0 & 2y(3) & 0 & 0 \\
0 & 0 & 2 & 0 & 0 & 0 & 0 & 2y(1) & 0 \\
0 & 0 & 0 & 6 & 2 & 6 & 0 & 0 & 2y(1) \\
0 & 0 & 0 & 6 & 6 & 2 & 0 & 0 & 2y(3) \\
0 & 0 & 0 & 0 & 0 & 0 & 2 & 2 & 0
\end{pmatrix}
$}
\]



\[
\chi_M(\lambda)
=
\lambda\Bigl(
\lambda^8
-16\bigl(y(1)+y(3)\bigr)\lambda^6
+\bigl(96y(1)^2-1984y(1)y(3)+96y(3)^2\bigr)\lambda^4
\]
\[
\hspace{3.5cm}
-\bigl(256y(1)^3+7936y(1)^2y(3)+7936y(1)y(3)^2+256y(3)^3\bigr)\lambda^2
+256\bigl(y(1)-y(3)\bigr)^4
\Bigr).
\]

Equivalently, setting \(t=\lambda^2\), the nonzero eigenvalues are obtained from
\[
t^4
-16\bigl(y(1)+y(3)\bigr)t^3
+\bigl(96y(1)^2-1984y(1)y(3)+96y(3)^2\bigr)t^2
\]
\[
\hspace{2.5cm}
-\bigl(256y(1)^3+7936y(1)^2y(3)+7936y(1)y(3)^2+256y(3)^3\bigr)t
+256\bigl(y(1)-y(3)\bigr)^4=0.
\]

So the eigenvalues are
\[
\lambda=0,
\qquad
\lambda=\pm\sqrt{t_1},\ \pm\sqrt{t_2},\ \pm\sqrt{t_3},\ \pm\sqrt{t_4},
\]
where \(t_1,t_2,t_3,t_4\) are the roots of the quartic above.


\subsection{$\mathcal{O}(1,1) \oplus \mathcal{O}(1,1)$}
The hodge diamond is
\[
    \begin{array}{ccc}
         &1&  \\
         &2&\\
         3&&3\\
         &2&\\
         &1&
    \end{array}
\]
% Prevent paragraph break confusion

\[
  \resizebox{\textwidth}{!}{$
\begin{pmatrix}
0 & 14y(1)y(3) & 14y(1)y(3) & 36y(1)^2y(3)+36y(1)y(3)^2 & 18y(1)^2y(3)+30y(1)y(3)^2 & 4y(1)^3y(3)+48y(1)^2y(3)^2+4y(1)y(3)^3 \\
1 & -y(1) & 3y(3) & 18y(1)y(3) & 10y(1)y(3) & 3y(1)^2y(3)+9y(1)y(3)^2 \\
1 & 3y(1) & -y(3) & 18y(1)y(3) & 14y(1)y(3) & 9y(1)^2y(3)+3y(1)y(3)^2 \\
0 & \tfrac{7}{3} & 1 & 3y(3)-y(1) & 4y(3) & \tfrac{14}{3}y(1)y(3) \\
0 & -1 & 1 & 3y(1)-3y(3) & -4y(3) & 0 \\
0 & 0 & 0 & 6 & 4 & 0
\end{pmatrix}
$}
\]


\[
\chi_M(\lambda)
=
(\lambda+y(1)+y(3))^2
\Bigl(
\lambda^4
-136\,y(1)y(3)\,\lambda^2
-384\,y(1)y(3)\bigl(y(1)+y(3)\bigr)\lambda
-16\,y(1)y(3)\bigl(16y(1)^2+31y(1)y(3)+16y(3)^2\bigr)
\Bigr).
\]

Hence the eigenvalues are
\[
\lambda_1=\lambda_2=-(y(1)+y(3)),
\]
and the remaining four eigenvalues are the roots of
\[
\lambda^4
-136\,y(1)y(3)\,\lambda^2
-384\,y(1)y(3)\bigl(y(1)+y(3)\bigr)\lambda
-16\,y(1)y(3)\bigl(16y(1)^2+31y(1)y(3)+16y(3)^2\bigr)=0.
\]
\subsection{$\mathcal{O}(2,2)$}

The Hoge diamond is
\[
\begin{array}{ccc}
     &1&  \\
     &2&\\
     16&114&16\\
     &2&\\
     &1&
\end{array}
\]
where $3$ out of the $114$ are ambient classes.


\section{OG(3,7)}

This is the quadric in $\mathbb{P}^7$.


\subsection{$\mathcal{O}(2)$}
Intersection of two quadrics is $\mathbb{P}^7$
\[
\begin{array}{ccc}
     &1&  \\
     &1&\\
     &1&\\
     3&&3\\
     &1&\\
     &1&\\
     &1&
\end{array}
\]
\[
\begin{pmatrix}
0 & 0 & 0 & 8y_{3} & 0 & 0 \\
4 & 0 & 0 & 0 & 16y_{3} & 0 \\
0 & 4 & 0 & 0 & 0 & 8y_{3} \\
0 & 0 & 8 & 0 & 0 & 0 \\
0 & 0 & 0 & 4 & 0 & 0 \\
0 & 0 & 0 & 0 & 4 & 0
\end{pmatrix}
\]

\[
\left\{
0,\,
0,\,
-8\sqrt[4]{y_{3}},\,
-8i\sqrt[4]{y_{3}},\,
8i\sqrt[4]{y_{3}},\,
8\sqrt[4]{y_{3}}
\right\}.
\]

\subsection{$\mathcal{O}(3)$}
Intersection of quadric and cubic in $\mathbb{P}^7$
\[
\begin{array}{ccccc}
     &&1&&  \\
     &&1&&\\
     &&1&&\\
     1&83&&83&1\\
        &&1&&  \\
     &&1&&\\
     &&1&&
\end{array}
\]
\[
\begin{pmatrix}
0 & 0 & 36y_{3} & 0 & 0 & 216y_{3}^{2} \\
3 & 0 & 0 & 63y_{3} & 0 & 0 \\
0 & 3 & 0 & 0 & 63y_{3} & 0 \\
0 & 0 & 6 & 0 & 0 & 36y_{3} \\
0 & 0 & 0 & 3 & 0 & 0 \\
0 & 0 & 0 & 0 & 3 & 0
\end{pmatrix}
\]
Eigenvalues $0,0,0,9 \cdot 2^{2/3}\xi^{1/3} , 9\cdot 2^{2/3}\xi^{2/3} , 9\cdot 2^{2/3}\xi$. Where $\xi$ is the 3-root of unity.


\subsection{$\mathcal{O}(4)$}
Intersection of quadric and quartic in $\mathbb{P}^7$
\[
\begin{array}{ccccc}
     &&1&&  \\
     &&1&&\\
     &&1&&\\
     36&783&&783&36\\
        &&1&&  \\
     &&1&&\\
     &&1&&
\end{array}
\]
\[
\begin{pmatrix}
0 & 96y_{3} & 0 & 26688y_{3}^{2} & 0 & 1603584y_{3}^{3} \\
2 & 0 & 512y_{3} & 0 & 78976y_{3}^{2} & 0 \\
0 & 2 & 0 & 416y_{3} & 0 & 26688y_{3}^{2} \\
0 & 0 & 4 & 0 & 512y_{3} & 0 \\
0 & 0 & 0 & 2 & 0 & 96y_{3} \\
0 & 0 & 0 & 0 & 2 & 0
\end{pmatrix}
\]
Eigenvalue
\[
\left\{
0,\,
0,\,
0,\,
0,\,
-64\sqrt{y_{3}},\,
64\sqrt{y_{3}}
\right\}.
\]
\subsection{$\mathcal{O}(2)\oplus \mathcal{O}(2)$}
Intersections of 3 quadrics in $\mathbb{P}^7$
\[
\begin{array}{ccc}
     &1&  \\
     &1&\\
     3&1+37&3\\
     &1&\\
     &1&
\end{array}
\]

\[
\begin{pmatrix}
0 & 16y_{3} & 0 & 224y_{3}^{2} & 0 \\
2 & 0 & 48y_{3} & 0 & 224y_{3}^{2} \\
0 & 2 & 0 & 24y_{3} & 0 \\
0 & 0 & 4 & 0 & 16y_{3} \\
0 & 0 & 0 & 2 & 0
\end{pmatrix}
\]

\[
\left\{
0,\,
0,\,
0,\,
-16\sqrt{y_{3}},\,
16\sqrt{y_{3}}
\right\}.
\]


\section{$LG(3,6)$}

\subsection{$\mathcal{O}(1)$}
\[
\begin{pmatrix}
0 & 0 & 6y_{3} & 0 & 0 & 0 \\
3 & 0 & 0 & 3y_{3} & 0 & 0 \\
0 & 6 & 0 & 0 & 6y_{3} & 0 \\
0 & 0 & 12 & 0 & 0 & 6y_{3} \\
0 & 0 & 0 & 3 & 0 & 0 \\
0 & 0 & 0 & 0 & 6 & 0
\end{pmatrix}
\]
\[
\left\{
3(-2)^{2/3}\big((3-2\sqrt{2})y_{3}\big)^{1/3},\,
3\cdot 2^{2/3}\big((3-2\sqrt{2})y_{3}\big)^{1/3},\,
-3(-1)^{1/3}2^{2/3}\big((3-2\sqrt{2})y_{3}\big)^{1/3},\right.
\]
\[
\left.
3(-2)^{2/3}\big((3+2\sqrt{2})y_{3}\big)^{1/3},\,
3\cdot 2^{2/3}\big((3+2\sqrt{2})y_{3}\big)^{1/3},\,
-3(-1)^{1/3}2^{2/3}\big((3+2\sqrt{2})y_{3}\big)^{1/3}
\right\}.
\]

\subsection{$\mathcal{O}(1)\oplus \mathcal{O}(1)$}
\[
\begin{pmatrix}
0 & 8y_{3} & 0 & 8y_{3}^{2} & 0 \\
2 & 0 & 8y_{3} & 0 & 8y_{3}^{2} \\
0 & 4 & 0 & 4y_{3} & 0 \\
0 & 0 & 8 & 0 & 8y_{3} \\
0 & 0 & 0 & 2 & 0
\end{pmatrix}
\]
\[
\left\{
0,\,
-4\sqrt{(3-2\sqrt{2})y_{3}},\,
4\sqrt{(3-2\sqrt{2})y_{3}},\,
-4\sqrt{(3+2\sqrt{2})y_{3}},\,
4\sqrt{(3+2\sqrt{2})y_{3}}
\right\}.
\]


\end{document}